\chapter{API y endpoints}

\section{Visión general}

La API expone el horario global, la información del grafo y los horarios personalizados de estudiantes. Todas las respuestas usan JSON y la capa de servicio mantiene en memoria el horario global ya calculado.

\section{GET /schedule}

\textbf{Request:}
\begin{lstlisting}[language=Bash]
GET /schedule
\end{lstlisting}

\textbf{Response:}
\begin{lstlisting}[language=JSON,caption={Respuesta de /schedule}]
{
  "success": true,
  "data": {
    "horario": [
      {
        "grupo": "SIS-ALG-G1",
        "materia": "Álgebra",
        "semestre": 1,
        "franja": "Lun-A1",
        "horario": "Lunes 08:00-10:00",
        "salon": "A101",
        "salonDetalle": {
          "id": "A101",
          "tipo": "normal",
          "capacidad": 40,
          "piso": 1,
          "sede": "Principal"
        },
        "tipoSalon": "normal",
        "capacidad": 40,
        "color": "A101"
      }
    ],
    "grafo": {
      "vertices": 108,
      "aristas": 1243
    }
  }
}
\end{lstlisting}

\section{GET /schedule/grafo}

\textbf{Request:}
\begin{lstlisting}[language=Bash]
GET /schedule/grafo
\end{lstlisting}

\textbf{Response:}
\begin{lstlisting}[language=JSON,caption={Respuesta de /schedule/grafo}]
{
  "success": true,
  "data": {
    "vertices": 108,
    "aristas": 1243
  }
}
\end{lstlisting}

\section{POST /student/schedule}

\textbf{Request:}
\begin{lstlisting}[language=JSON,caption={Body de /student/schedule}]
{
  "estudianteId": "EST-20221020084",
  "materias": ["SIS-ALG", "SIS-BD1", "SIS-PROG"]
}
\end{lstlisting}

\textbf{Response:}
\begin{lstlisting}[language=JSON,caption={Respuesta de /student/schedule}]
{
  "success": true,
  "data": {
    "estudianteId": "EST-20221020084",
    "bloques": [
      {
        "materiaId": "SIS-ALG",
        "nombreMateria": "Álgebra",
        "grupo": "SIS-ALG-G1",
        "dia": "Lunes",
        "horaInicio": "08:00",
        "horaFin": "10:00",
        "salon": "A101",
        "sede": "Principal",
        "piso": 1,
        "tipoSalon": "normal",
        "bloque": 1
      }
    ],
    "metricas": {
      "totalBloques": 3,
      "cambiosDeSede": 1,
      "costoTotal": 4,
      "sedesUsadas": ["Principal"],
      "materiasNoAsignadas": []
    }
  }
}
\end{lstlisting}

\section{GET /student/schedule/:estudianteId}

\textbf{Request:}
\begin{lstlisting}[language=Bash]
GET /student/schedule/EST-20221020084
\end{lstlisting}

\textbf{Response:}
\begin{lstlisting}[language=JSON,caption={Respuesta de /student/schedule/:estudianteId}]
{
  "success": true,
  "data": {
    "estudianteId": "EST-20221020084",
    "bloques": [/* ... */],
    "metricas": { /* ... */ }
  }
}
\end{lstlisting}

\section{GET /student/schedules}

\textbf{Request:}
\begin{lstlisting}[language=Bash]
GET /student/schedules
\end{lstlisting}

\textbf{Response:}
\begin{lstlisting}[language=JSON,caption={Respuesta de /student/schedules}]
{
  "success": true,
  "data": {
    "total": 1,
    "estudiantes": [
      {
        "estudianteId": "EST-20221020084",
        "totalBloques": 3,
        "costoTotal": 4,
        "cambiosDeSede": 1,
        "sedesUsadas": ["Principal"],
        "materiasNoAsignadas": []
      }
    ]
  }
}
\end{lstlisting}

\section{Formato de respuesta}

Todas las rutas siguen la estructura:
\begin{lstlisting}[language=TypeScript]
interface ApiResponse<T> { success: boolean; data?: T; error?: string }
\end{lstlisting}
